\documentclass[12pt]{article}

\usepackage[margin=1in]{geometry}
\usepackage{amsmath,amssymb,amsthm}
\usepackage{natbib}
\usepackage[hidelinks]{hyperref}

\pagestyle{empty}
\newtheorem{proposition}{Proposition}
\newtheorem{corollary}{Corollary}
\newtheorem{lemma}{Lemma}

\makeatletter
\renewcommand{\@seccntformat}[1]{\csname the#1\endcsname.\quad}
\renewcommand\section{\@startsection{section}{1}{\z@}%
  {-2.5ex \@plus -1ex \@minus -.2ex}%
  {1.3ex \@plus .2ex}%
  {\centering\normalfont\fontsize{14}{16}\selectfont\bfseries}}
\makeatother

\begin{document}

\section{Introduction}

Regional industrial policy is often discussed in terms of how much support governments should provide. A distinct short-run problem arises once the policy envelope is predetermined and fully committed: where should that envelope be directed? Regional governments may favor the same attractive industries, yet simultaneous targeting can crowd the same projects, specialized workers, suppliers, infrastructure, or administrative resources. The resulting problem concerns policy composition rather than policy level. This paper asks whether decentralized regional competition can sustain duplicated sectoral priorities even when constrained coordination would favor differentiation.

The question connects two established literatures. First, fiscal-competition models show that decentralization can distort the composition of public spending. \citet{keen_marchand_1997} establish a bias between productive public inputs and public goods for immobile residents; subsequent work studies the robustness of that composition effect to factor mobility and richer spending categories \citep{matsumoto_2000,borck_2005,borck_caliendo_steiner_2007}. \citet{arcalean_et_al_2010} likewise study budget composition across education and infrastructure in a federal economy in which both budget size and composition are endogenous. Second, the public-input competition literature shows that regional competition can generate inefficient spatial patterns. \citet{bucovetsky_2005} demonstrates that too many regions may invest in public inputs, while \citet{fenge_ehrlich_wrede_2009} show that decentralized public-input competition and the efficient allocation can imply different spatial concentration patterns. These literatures establish broad possibilities of composition distortions and inefficient spatial allocation; the contribution here is narrower.

Each of two identical regional governments has a predetermined, fully committed policy envelope $B$ and allocates it between two productive sectors. Support withdrawn from one sector must be reassigned one-for-one to the other. This equality constraint couples two sectoral contests that would otherwise be independent. Sector $A$ has a common intrinsic advantage $\Delta>0$, own policy has diminishing returns $c>0$, and simultaneous same-sector targeting generates a real cross-regional overlap loss $\rho\ge0$.

The main result contrasts decentralized duplication with constrained differentiation. For $0\le\rho<c$, the decentralized policy game has a unique Nash equilibrium, and both governments are oriented toward the same intrinsically more attractive sector $A$. Yet a coordinator subject to the same region-specific envelopes assigns opposite sectoral orientations if and only if
\[
\rho>\frac{c}{2}
\qquad\text{and}\qquad
B>\frac{\Delta}{2\rho}.
\]
Thus rivalry can be strong enough for constrained coordination to favor differentiation without being strong enough to induce decentralized differentiation. The envelope threshold has a direct interpretation: as the committed envelope expands, the feasible scope for differentiation, and hence the overlap loss that coordination can avoid, rises relative to the fixed intrinsic advantage of sector $A$.

The equality constraint is essential to this $B$-driven result. If sectoral policy levels are chosen independently, the payoff separates into two standard policy-level games and $B$ disappears, so there is no analogous threshold. The claim is deliberately narrow: the paper isolates a composition mechanism created by competition over the sectoral allocation of a predetermined, fully committed policy envelope.

\section{Model}

Consider two identical regions, indexed by $i\in\{1,2\}$, and two productive sectors, $A$ and $B$. Each regional government has a predetermined and fully committed short-run industrial-policy envelope $B>0$. Region $i$ allocates $x_i\in[0,B]$ to sector $A$ and the residual $B-x_i$ to sector $B$. The equality constraint fixes the total committed policy quantity, so only sectoral composition is strategic; the prior decision of how much policy to activate lies outside the model.

Sector $A$ has a common intrinsic advantage $\Delta>0$ over sector $B$. Own-region policy exhibits diminishing returns, governed by $c>0$, while simultaneous targeting of the same sector by the two regions creates a real overlap cost $\rho\ge 0$. Region $i$'s payoff is
\begin{equation}
\begin{split}
U_i(x_i,x_j)= {} &(a+\Delta)x_i-\frac{c}{2}x_i^2-\rho x_i x_j \\
&+a(B-x_i)-\frac{c}{2}(B-x_i)^2
-\rho(B-x_i)(B-x_j),
\end{split}
\label{eq:payoff}
\end{equation}
where $j\neq i$. Throughout the main analysis,
\begin{equation}
\Delta>0,\qquad c>0,\qquad 0\le \rho<c,\qquad B>0.
\label{eq:domain}
\end{equation}
The restriction $\rho<c$ isolates the region in which decentralized portfolio choice has a unique equilibrium. The stronger-rivalry case is not needed for the result below.

A simple project-pool interpretation delivers \eqref{eq:payoff} exactly. Suppose sector $s$ contains a mass $M$ of potential projects and that targeting intensity $z_{is}\in[0,B]$ reaches any project with probability $q_{is}=z_{is}/\bar z$, where $\bar z\ge B$. A project targeted by region $i$ yields local value $v_s$, while overlapping targeting by both regions generates a real implementation or congestion loss $\ell>0$ for each region. With convex implementation cost $\kappa z_{is}^2/2$, expected sectoral payoff is
\[
M\left(v_s q_{is}-\ell q_{is}q_{js}\right)-\frac{\kappa}{2}z_{is}^2
=a_s z_{is}-\frac{c}{2}z_{is}^2-\rho z_{is}z_{js},
\]
where $a_s=Mv_s/\bar z$, $c=\kappa$, and $\rho=M\ell/\bar z^2$. Setting $a_A=a+\Delta$ and $a_B=a$ gives \eqref{eq:payoff}. Thus $\rho$ captures a real same-sector overlap loss rather than a pure transfer of subsidy rents. This real-resource interpretation gives the coordination comparison its normative content within the reduced-form environment.

We call a region \emph{$A$-oriented} if $x_i>B/2$, \emph{$B$-oriented} if $x_i<B/2$, and \emph{neutral} if $x_i=B/2$. The two regions display \emph{policy duplication} when they have the same sectoral orientation. They display \emph{strict sectoral differentiation} when, up to relabeling, $x_1>B/2>x_2$.

\section{Decentralized Portfolio Choice}

Differentiating \eqref{eq:payoff} gives
\begin{equation}
\frac{\partial U_i}{\partial x_i}
=\Delta+B(c+\rho)-2cx_i-2\rho x_j,
\qquad
\frac{\partial^2 U_i}{\partial x_i^2}=-2c<0.
\label{eq:foc}
\end{equation}
Hence the global best response is
\begin{equation}
BR_i(x_j)=
\Pi_{[0,B]}
\left[
\frac{\Delta+B(c+\rho)-2\rho x_j}{2c}
\right],
\label{eq:br}
\end{equation}
where $\Pi_{[0,B]}$ denotes projection onto $[0,B]$.

\begin{proposition}[Decentralized portfolio choice]
Under \eqref{eq:domain}, the policy game has a unique Nash equilibrium. It is
\begin{equation}
(x_1^N,x_2^N)=
\begin{cases}
(B,B), & B\le \dfrac{\Delta}{c+\rho},\\[6pt]
\left(
\dfrac{B}{2}+\dfrac{\Delta}{2(c+\rho)},
\dfrac{B}{2}+\dfrac{\Delta}{2(c+\rho)}
\right),
& B>\dfrac{\Delta}{c+\rho}.
\end{cases}
\label{eq:ne}
\end{equation}
In particular, both regions are $A$-oriented for every admissible parameter vector.
\end{proposition}

\begin{proof}
Because projection is nonexpansive, \eqref{eq:br} implies
\[
|BR_i(x)-BR_i(y)|\le \frac{\rho}{c}|x-y|.
\]
Hence the joint best-response mapping is a contraction in the sup norm with modulus $\rho/c<1$, so it has a unique fixed point. Because the game is symmetric, swapping the coordinates of any fixed point produces another fixed point; uniqueness therefore implies $x_1^N=x_2^N=x^N$. If the upper bound is slack, the fixed-point equation from \eqref{eq:br} gives
\[
x^N=\frac{B}{2}+\frac{\Delta}{2(c+\rho)}.
\]
This value is below $B$ exactly when $B>\Delta/(c+\rho)$. Otherwise the unique fixed point is $(B,B)$. In both cases $x_i^N>B/2$ because $\Delta>0$.
\end{proof}

Strategic substitutability alone does not generate different regional sectoral priorities. Even though $BR_i'(x_j)=-\rho/c<0$ wherever the best response is interior, both governments continue to favor the same intrinsically more attractive sector.

\section{Constrained Coordination}

To isolate the portfolio externality, consider a coordinator who maximizes the aggregate regional payoff represented by the model,
\[
W(x_1,x_2)=U_1(x_1,x_2)+U_2(x_2,x_1),
\]
subject to the same region-specific envelope constraints $x_i\in[0,B]$. The coordinator cannot change either region's total committed envelope or transfer policy capacity across regions; it only coordinates sectoral composition. Thus $W$ is the constrained-coordination objective within this reduced-form environment, not an unrestricted national social-welfare function.

Define
\begin{equation}
p=\frac{x_1+x_2}{2}-\frac{B}{2},
\qquad
d=\frac{x_1-x_2}{2}.
\label{eq:pd}
\end{equation}
Here $p$ measures the common tilt toward sector $A$ relative to an even split, while $d$ measures cross-regional dispersion. The feasible set becomes $|p|+|d|\le B/2$. Up to an additive constant independent of $(p,d)$,
\begin{equation}
W(p,d)=2\Delta p-2(c+2\rho)p^2-2(c-2\rho)d^2.
\label{eq:welfare}
\end{equation}
The coefficient on $d^2$ separates the decentralized and coordinated incentives. When $\rho<c/2$, the coordinator favors equal regional portfolios. When $\rho>c/2$, holding the mean portfolio fixed, the coordinator objective rises as the two regional portfolios move apart.

\begin{proposition}[Constrained coordination]
Suppose \eqref{eq:domain} holds.
\begin{enumerate}
\item If $0\le\rho<c/2$, every coordinated optimum is symmetric. It equals $(B,B)$ when $B\le \Delta/(c+2\rho)$ and otherwise satisfies
\[
x_1^P=x_2^P=\frac{B}{2}+\frac{\Delta}{2(c+2\rho)}.
\]
\item If $\rho>c/2$, the coordinated optimum is unique up to regional relabeling. For $B\le \Delta/(c+2\rho)$ it is $(B,B)$. For larger $B$ it is
\begin{equation}
(B,z^P)\quad\text{or}\quad(z^P,B),
\qquad
z^P=\max\left\{0,\frac{\Delta+B(c-2\rho)}{2c}\right\}.
\label{eq:planner}
\end{equation}
\end{enumerate}
At the knife edge $\rho=c/2$, $(B,B)$ is the unique coordinated optimum if $B\le\Delta/(2c)$. If $B>\Delta/(2c)$, the coordinated optima are exactly the feasible allocations satisfying
\[
x_1+x_2=B+\frac{\Delta}{2c}.
\]
\end{proposition}

\begin{proof}
If $\rho<c/2$, \eqref{eq:welfare} is strictly decreasing in $d^2$, so every optimum has $d=0$. The remaining concave quadratic in $p$ has unconstrained maximizer $\Delta/[2(c+2\rho)]$. Projecting this value onto $[-B/2,B/2]$ gives the stated symmetric solution, with the upper boundary reached exactly when $B\le\Delta/(c+2\rho)$.

If $\rho>c/2$, \eqref{eq:welfare} is increasing in $d^2$ for fixed $p$, so an optimum satisfies $|d|=B/2-|p|$. For any $p>0$, the feasible points with common tilts $p$ and $-p$ have the same quadratic terms in \eqref{eq:welfare}, while their objective values differ by $4\Delta p>0$. Hence an optimum has $p\ge0$. Up to regional relabeling, set $d=B/2-p$, which gives $(x_1,x_2)=(B,2p)$. Writing $z=2p\in[0,B]$ and substituting into $W$ yields a strictly concave quadratic in $z$ whose unconstrained maximizer is
\[
\frac{\Delta+B(c-2\rho)}{2c}.
\]
Clipping this value to $[0,B]$ gives \eqref{eq:planner}. The upper bound binds exactly when $B\le\Delta/(c+2\rho)$.

At $\rho=c/2$, the $d^2$ term in \eqref{eq:welfare} vanishes. Maximizing over $p$ gives $p^P=\min\{B/2,\Delta/(4c)\}$. If $B\le\Delta/(2c)$, then $p^P=B/2$ and feasibility forces $d=0$, yielding $(B,B)$. If $B>\Delta/(2c)$, then $p^P=\Delta/(4c)$ and every feasible $d$ is optimal, which is equivalent to $x_1+x_2=B+\Delta/(2c)$.
\end{proof}

For $\rho>c/2$, asymmetry begins once
\[
B>\frac{\Delta}{c+2\rho}.
\]
This is not yet sectoral differentiation: for an intermediate range both regions remain $A$-oriented. Strict differentiation requires $z^P<B/2$, which is equivalent to
\begin{equation}
B>\frac{\Delta}{2\rho}.
\label{eq:diffthreshold}
\end{equation}
At equality, the low-$A$ region is neutral, $z^P=B/2$.

\section{Excessive Policy Duplication}

\begin{corollary}[Excessive policy duplication]
Under \eqref{eq:domain}, the unique decentralized equilibrium is $A$-duplicated and every constrained coordinated optimum is strictly sectorally differentiated if and only if
\begin{equation}
\rho>\frac{c}{2}
\qquad\text{and}\qquad
B>\frac{\Delta}{2\rho}.
\label{eq:headline}
\end{equation}
Equivalently, in the headline region $c/2<\rho<c$ and $B>\Delta/(2\rho)$, both decentralized governments are $A$-oriented, whereas every coordinated optimum has one $A$-oriented region and one $B$-oriented region.
\end{corollary}

\begin{proof}
Proposition 1 implies that the decentralized equilibrium is unique and $A$-duplicated throughout the production domain. Proposition 2 rules out strict coordinated differentiation when $\rho<c/2$, while at $\rho=c/2$ a symmetric coordinated optimum remains available. Now suppose $\rho>c/2$. If $B\le\Delta/(c+2\rho)$, the coordinated optimum is $(B,B)$. For larger $B$, \eqref{eq:planner} is strictly differentiated exactly when $z^P<B/2$. When $z^P>0$,
\[
z^P-\frac{B}{2}=\frac{\Delta-2\rho B}{2c},
\]
so $z^P<B/2$ if and only if $B>\Delta/(2\rho)$. If the lower bound binds, then $\Delta+B(c-2\rho)\le0$, so $B\ge\Delta/(2\rho-c)>\Delta/(2\rho)$ and strict differentiation again holds. Thus \eqref{eq:headline} is necessary and sufficient.
\end{proof}

The result isolates an intermediate rivalry region. Rivalry can be strong enough for constrained coordination to assign opposite sectoral orientations, yet too weak to break the unique decentralized symmetry. The committed envelope then determines whether this coordination motive overturns the common intrinsic advantage of sector $A$. Rearranging \eqref{eq:diffthreshold}, the constrained optimum is strictly differentiated when
\[
\rho B>\frac{\Delta}{2}.
\]
A larger committed envelope therefore expands the feasible scope for differentiation and the overlap loss that coordination can avoid, even though each region's total committed policy quantity remains fixed.

The wedge reflects a portfolio-overlap externality. A marginal increase in region $i$'s allocation to sector $A$ changes region $j$'s payoff by
\begin{equation}
\frac{\partial U_j}{\partial x_i}=\rho(B-2x_j).
\label{eq:externality}
\end{equation}
Whenever region $j$ is $A$-oriented, this effect is negative. Region $i$ internalizes how region $j$'s targeting affects its own return, but not the reciprocal effect \eqref{eq:externality} that its portfolio imposes on the other region. Constrained coordination internalizes both sides.

\section{Why the Capacity Constraint Matters}

The fixed-envelope equality constraint is not a normalization; it creates the portfolio problem. To see this, remove the constraint and allow each government to choose $z_{iA}\ge0$ and $z_{iB}\ge0$ independently. Then
\begin{equation}
U_i=\sum_{s\in\{A,B\}}
\left(a_s z_{is}-\frac{c}{2}z_{is}^2-\rho z_{is}z_{js}\right).
\label{eq:unconstrained}
\end{equation}
The two sectors are additively separable, with best response
\[
BR_{is}(z_{js})=
\max\left\{0,\frac{a_s-\rho z_{js}}{c}\right\}.
\]
Each sector can therefore be solved as an independent policy-level game. The envelope parameter $B$ disappears from both the strategy set and the payoff, so there is no threshold analogous to \eqref{eq:diffthreshold}. Under the fixed-envelope equality constraint, by contrast, withdrawing one unit of support from one sector necessarily reallocates it to the other. This one-for-one coupling generates the $B$-driven duplication result. If $B$ were instead only an optional ceiling with unused capacity permitted, total policy quantity would again become a choice margin, defining a different model.

\section{Conclusion}

With a predetermined, fully committed short-run policy envelope, decentralized competition can preserve duplicated sectoral priorities even when constrained coordination favors differentiation. In the benchmark model, this occurs exactly when rivalry exceeds $c/2$ and the committed envelope exceeds $\Delta/(2\rho)$. The mechanism is portfolio coupling: support withdrawn from one sector must be reassigned to the other.

The result is deliberately conditional. It does not imply that industrial-policy budgets should be centralized, that duplication is always inefficient, or that larger policy budgets are undesirable. The coordinator cannot change or transfer regional policy envelopes; it only internalizes the portfolio-overlap externality. The implication is narrower: a larger committed envelope can move the constrained-coordination problem from duplication toward differentiation even when the prior decision over how much policy to activate lies outside the model. Removing the fixed-envelope equality constraint eliminates this comparative static, identifying portfolio coupling as the mechanism behind the result.

\bibliographystyle{apalike}
\bibliography{references}

\end{document}
